\documentclass[12pt]{report}

\usepackage{amsmath, amssymb, amsfonts}
\usepackage{geometry}
\usepackage{setspace}
\usepackage{titlesec}
\usepackage{hyperref}

\geometry{margin=1in}
\setstretch{1.2}

\titleformat{\chapter}{\normalfont\Large\bfseries}{\thechapter}{1em}{}
\titleformat{\section}{\normalfont\large\bfseries}{\thesection}{1em}{}
\titleformat{\subsection}{\normalfont\normalsize\bfseries}{\thesubsection}{1em}{}

\begin{document}

\begin{titlepage}
    \centering
    {\Large\bfseries SEFI/GWFM: Unified Geometric Field Theory}\\

    \vspace{1.5em}
    {\large Jason Duran Dutton}\\

    \vspace{0.5em}
    {\large Castlewood, Virginia, USA}\\

    \vspace{2em}
    {\large 2026}
\end{titlepage}

\tableofcontents
\newpage

%===========================================================
\chapter{Worldline Foundations (GWFM)}
%===========================================================

\section{Introduction}

Quantum mechanics and quantum field theory accurately predict experimental
outcomes but leave open the ontological status of particles, fields, and
superposition. The Geometric Worldline Field Model (GWFM) begins with a
geometric premise: a continuous worldline entity spans spacetime, and the
electron field is the projection of this structure. Superposition, interference,
and measurement arise from the geometry of this worldline rather than from
separate interpretive rules.

\section{Continuous Worldline Entity}

GWFM posits that the electron field arises from a single continuous worldline
entity embedded in spacetime. Let $M$ be a spacetime manifold equipped with
metric $g_{\mu\nu}$. The entity's worldline is defined as a continuous map


\[
    y : \mathbb{R} \rightarrow M,
\]


where the parameter $T$ indexes the full geometric extent of the entity. The
worldline spans all allowed paths consistent with the worldline action, producing
a dense geometric ensemble whose projection into spacetime yields the appearance
of a continuous electron field.

Forward-time segments satisfy $\frac{dt}{dT} > 0$ (electrons), while
backward-time segments satisfy $\frac{dt}{dT} < 0$ (positrons). This orientation
rule provides a geometric interpretation of particle--antiparticle duality.

Superposition arises from the entity spanning multiple geometric branches,
interference arises from overlap regions of the worldline, and measurement
corresponds to boundary-condition selection of specific segments.

\section{Dense Trajectory Ensemble}

The worldline spans all allowed trajectories, forming a dense geometric set.
Presence-density is defined as


\[
p(x) = \int Dy\,D\xi\, e^{S_{\text{entity}}[y,\xi]}\, \delta(x - y(T)).
\]



\section{Field Emergence}

The electron field is the projection of the worldline ensemble:


\[
\phi(x) = \int dT\, \delta(x - y(T)).
\]



\section{Spin Structure}

Spin degrees of freedom are encoded using Grassmann variables $\xi^\mu(T)$ with
action


\[
S_{\text{spin}} = \int dT\, \xi_\mu \dot{\xi}^\mu.
\]



\section{Gauge Coupling}

Electromagnetic coupling is introduced through


\[
S_{\text{EM}} = q \int dT\, \dot{y}^\mu A_\mu(y(T)).
\]



\section{Interpretation}

Interference corresponds to geometric overlap, measurement corresponds to
boundary selection, and identity is global and geometric.

%===========================================================
\chapter{Layered Field Identity (SEFI)}
%===========================================================

\section{Introduction}

SEFI extends GWFM by introducing layered field identity: FIELD ORIGIN, FIELD
AUTHORSHIP, FIELD SOVEREIGNTY, WARP EXPRESSION, and the R-LAYER (Reality
Field). These layers define stability, identity, and transformation rules.

\section{FIELD ORIGIN}

FIELD ORIGIN defines primitive geometric identity. Stability is treated as a
primitive property:


\[
\left\| \frac{d^2 r}{dt^2} \right\| \le S_0,
\]


where $S_0$ is the sovereignty-bound acceleration limit.

\section{FIELD AUTHORSHIP}

Authorship layers contribute geometric identity constructs:


\[
F = \sum_{i=1}^N A_i(r(t), K(t)),
\]


where $K(t)$ is curvature.

\section{FIELD SOVEREIGNTY}

Sovereignty ensures invariance under allowed transformations:


\[
F(T(r(t))) = F(r(t)).
\]



\section{WARP EXPRESSION}

Warp operators modify worldline geometry while preserving identity:


\[
r'(t) = W(r(t)), \qquad F' = F(W(r(t))).
\]



\section{Reality Field Layer (R-Layer)}

The R-layer defines the actual geometric manifestation of the entity's worldline
into spacetime. It maps worldline geometry, curvature, torsion, stability
constraints, sovereignty invariants, and warp transformations into the observable
field.

\section{Mathematical Formulation}

Curvature is defined as


\[
K(t) = \frac{\| r'(t) \times r''(t) \|}{\| r'(t) \|^3}.
\]



The full SEFI identity is


\[
F(t) = O(r(t)) + A(r(t), K(t)) + S(r(t)) + W(r(t)) + R(r(t)).
\]



%===========================================================
\chapter{Photonic QEC from SEFI/GWFM Geometry}
%===========================================================

\section{Introduction}

Photonic systems suffer from loss, dephasing, mode mixing, and spectral drift.
SEFI/GWFM provides a geometric foundation for photonic QEC, defining stability
surfaces, geometric stabilizers, and syndrome geometry.

\section{Stability Surfaces}

Sovereignty defines a stability surface: a geometric manifold representing allowed
transformations. Errors correspond to displacements orthogonal to this surface.

\section{Photonic Error Geometry}

Loss corresponds to worldline termination, dephasing to torsion, mode mixing to
displacement into neighboring surfaces, and spectral wandering to translation
along the spectral axis.

\section{Geometric Stabilizers}

Stabilizers preserve sovereignty through curvature, torsion, spectral, and
polarization constraints.

\section{Syndrome Geometry}

Weak or QND measurements sample curvature deviation, torsion deviation, spectral
displacement, and polarization rotation.

\section{SPDC Geometry}

SPDC annular geometry corresponds to a stability surface defined by FIELD ORIGIN
and FIELD SOVEREIGNTY. Diametric correlations arise from sovereign geometric
symmetry.

\section{CHSH Geometry}

Polarization evolution is treated as a worldline on the Poincaré-Bloch sphere.
Bell states correspond to minimal torsion worldlines.

%===========================================================
\chapter{Conclusion}

SEFI/GWFM provides a unified geometric field theory. GWFM defines worldline
foundations, SEFI defines layered identity, and photonic QEC geometry provides
applications. Together they form a coherent geometric interpretation of field
behavior, stability, and transformation.

\end{document}
